\scanpage{95}
\mspage{47}
qui nous incite justement : regarder la \unsure{classes|classe} de scindages, de $\mathcal E_{X_{\unsure{B.|B_0}}}$ sur $\mathcal E_{k(s)}$ relatifs à cet $L_i$ --- et notamment ceux relatifs au choix \unsure{d'un|d'une} uniformisant de $\mathcal O_{S,s_\alpha}$, i.e. d'une équation locale de $D_\alpha$ au vois. de son pt gén. $s_\alpha$, \deleted{\unsure{(plus gén.)}} \inserted{(qui \unsure{intervient|interviennent} \unsure{p\textsuperscript{r}|par} le \unsure{brin|biais|bon|bout} qu'elles définissent $\mathfrak m_{s_\alpha}/\mathfrak m_{s_\alpha}^{2}$)} Quant à ces sections-là, on voit comment elles se spécialisent, du moins c'est clair en les pts de $D_\alpha$ \unsure{sur|en} les quels la section rationnelle envisagée du faisceau conormal $\underline N_{D_\alpha,S}^{\unsure{r|+|\vee}}$ (si on suppose $D_\alpha$ régulier, \unsure{\deleted{)}|)|\&} \unsure{de|donc|du} \unsure{..|:} $\underline N_{D_\alpha/S}$ inversible) reste régulière, et inversible.

On peut se proposer d'étudier \deleted{comment} \inserted{ce qui} se passe en les autres points --- i.e. la spécialisation des sections \inserted{de deuxième type} d'extensions de groupes profinis, en faisant une étude analogue à la précédente, pour cette question, Mais cela me semble \deleted{\unsure{unique|moins}} \inserted{d'un} intérêt moins brulant --- car cela \unsure{vient|n'est}
\margin{94}
